\section{Motivation and Requirements}

\subsection{Natural-Language Harnesses Are Real}

We built a snapshot and analysis pipeline for
\texttt{CLAUDE.md} and \texttt{AGENTS.md} files from popular software projects.
The preliminary extraction shows that these files contain many imperative rules,
but only a subset are OS-observable behavioral constraints. This distinction is
important. A rule such as ``never use \texttt{any}'' is a coding-style rule and
is outside ActPlane's scope. A rule such as ``never push directly to
\texttt{main}'' or ``run tests before committing'' has an observable process
and syscall footprint and can be enforced as a harness contract.

The paper will use the corpus study for two purposes. First, it grounds the
motivation: developers already write behavioral constraints for agents. Second,
it derives evaluation workloads: only rules that are observable, expressible,
robust, and temptable are converted into policy/task pairs. The final paper
should report both the unbiased corpus denominator and the filtered evaluation
subset, rather than selecting only examples that ActPlane can handle.

\subsection{Why Tool-Layer Guards Are Insufficient}

Agent frameworks and permission systems can intercept tool calls and apply
deterministic policies~\cite{wang2025agentspec,shi2025progent}. This is useful,
but it assumes that effects pass through the guarded tool boundary. A coding
agent violates that assumption whenever it invokes a shell, starts a subprocess,
loads a library, or talks to an external service from generated code. The same
semantic action, for example invoking \texttt{git push --force}, can appear as a
direct shell command, a nested \texttt{bash -c} string, a Python subprocess, or
a helper binary. A tool-layer guard sees only the outer tool invocation.

The OS sees the effect itself. Every exec, file open, unlink, rename, truncate,
and socket connect crosses a kernel boundary. If a policy is phrased over those
effects and over the lineage of the agent process tree, then the policy applies
regardless of the high-level route the agent selected.

\subsection{Requirements}

ActPlane is designed around five requirements.

\textbf{R1: Whole-execution scope.} A rule must apply to the agent's
descendants, not just to the top-level process or the initial tool call.

\textbf{R2: Cross-object provenance.} Many rules are about information flow,
not a single event. Reading a secret file should taint the reader; writing a
derived file should taint that file; reading the derived file later should taint
the next process.

\textbf{R3: Effect semantics that match the backend.} \texttt{block} means
pre-operation denial and is implemented only by BPF-LSM hooks. Tracepoints do
not support \texttt{block}; they can report \texttt{audit} rules or terminate
\texttt{kill} rules.

\textbf{R4: Corrective feedback.} The harness should not only fail a violating
action. It should tell the agent which rule fired, why, and what alternative or
precondition can allow progress.

\textbf{R5: Deployable mechanism.} The implementation should use upstream Linux
interfaces: eBPF for dynamic loading and BPF-LSM for pre-operation deny, rather
than a custom kernel module.
